\documentclass[10pt]{article}
\usepackage[margin=20mm]{geometry}
\usepackage{amsmath,amssymb,mathtools,fontspec,unicode-math,xurl,hyperref}
\IfFontExistsTF{DejaVu Serif}{\setmainfont{DejaVu Serif}}{\setmainfont{Latin Modern Roman}}
\setmathfont{latinmodern-math.otf}
\setlength{\parindent}{0pt}\setlength{\parskip}{6pt}\sloppy
\title{Scientific initial review of
The eighth moment of the Riemann zeta function}
\author{AI-assisted technical review}
\date{26 October 2026}
\begin{document}
\maketitle
Abstract: One local sign error is established in the second-order Euler-factor expansion in Proposition 6.5(ii). The exact factor, the proof later in the article, and an admissible asymptotic specialization agree on the opposite sign. The demonstrated impact is local: the required convergence mechanism and the leading eighth-moment coefficient survive this correction. The inspected conditional main-theorem mechanism yields no further established substantive issue. This initial assessment does not certify the entire paper; the scientific decision remains pending.\par
\section*{Identity, hypotheses, and scope}
The reviewed work is by Nathan Ng, Quanli Shen, and Peng-Jie Wong, DOI https://doi.org/10.4171/JEMS/1568. The exact reviewed publisher PDF is Journal of the European Mathematical Society 28 (2026), no. 10, 4175–4237, published online on 9 January 2025. It has 63 pages. Printed page numbers are used below, with PDF page \(j\) corresponding to printed page \(4174+j\). The separate author source is arXiv:2204.13891v1, submitted 29 April 2022; its identity and text are not treated as identical to the publisher version.\par
Theorem 1.1, p. 4178 (PDF 4), assumes both the Riemann hypothesis and Conjecture 1, the smoothed quaternary additive divisor conjecture. Its conclusion is the leading asymptotic \[I_4(T)=\int_0^T|\zeta(\tfrac12+it)|^8\,dt
\sim \frac{24024a_4}{16!}\,T(\log T)^{16}.\] It does not assert an unconditional theorem or a complete polynomial of lower-order terms. Conjecture 1 fixes a constant \(C>0\), admits small positive \(\varepsilon_0,\varepsilon'\), and requires the additive shift to satisfy \(1\le |r|\ll X^{1-\varepsilon'}\); its error is \(O(P^CX^{1/2+\varepsilon_0})\). These dependencies were retained in the checks.\par
Actual review concentrated on the statements and normalization, truncation and removal of smoothing, the application of the divisor conjecture, selected contour and Euler-product steps, and the leading residue arithmetic. The source inventory was also mechanically verified in full. That operation is not a scientific check. The report does not certify every derivative estimate, high-order remainder, external theorem, or argument on all 63 pages.\par
\section*{Established problem}
Problem 1 (Minor: the sign of the second Euler-factor coefficient).  Location. Proposition 6.5(ii), p. 4211 (PDF 37), with \(Y_p\) defined in (6.30); the remainder exponent on p. 4212 (PDF 38); the exact factor (6.29); and the proof on pp. 4227–4228 (PDF 53–54).\par
Equation reference: eq:printed. Claim. The proposition states \[
\mathcal I_p(u_1,u_2,u_3;s)
=1-Y_p p^{-2-4s}+O\!\left(p^{7\delta'+\vartheta(\sigma)}\right),
\qquad \sigma=\Re s,\] for \(|u_i|<\delta'\), \(0<\delta'<1/5\), and \(\sigma\ge-1/2+2\delta'\). In particular, \(\vartheta(\sigma)=-2-2\sigma\) when \(-1/4\le\sigma<0\).\par
Decisive coefficient calculation. Set \(a=X_1^{-1}\), \(b=X_2^{-1}\), and \(W=p^{-1-2s}\). The leading arithmetic coefficients in the exact factor, as expanded in the proof, are \[A_0=(4-a)(4-b),\qquad B_0=(10-4a)(10-4b).\] The normalization factor is \[\Pi=(1-W)^{16}(1-abW)(1-aW)^{-4}(1-bW)^{-4}
=1+a_1W+a_2W^2+\cdots,\] where \[\begin{aligned}
a_1&=-A_0,\\
a_2&=120+32ab+10(a^2+b^2)-64(a+b)-4(a^2b+ab^2).\end{aligned}\] Multiplication shows that the coefficient of \(W^2\) is \[\begin{aligned}
B_0+A_0a_1+a_2
&=-6a^2-6b^2-16ab+4a^2b+4ab^2-a^2b^2\\
&\hspace{1.5em}-36+24a+24b
=Y_p .\end{aligned}\] Thus the sign is \(+Y_pW^2\). The proof on p. 4228 itself gives this same positive coefficient before identifying it with \(Y_p\). This is an internal disagreement between the statement and the demonstrated expansion.\par
Why the discrepancy is not hidden by the error term. Take \[u_1=u_2=u_3=0,\qquad \delta'=1/100,\qquad s=-1/10.\] All displayed assumptions hold. Here \(a=b=1\), \(Y_p=-9\), and \(W=p^{-4/5}\). From Lemma 2.3 and (6.29), \[G_4(1;p)=3-\frac3p+\frac1{p^2},\qquad
G_4(1;p^2)=6-\frac8p+\frac3{p^2},\] and \[\mathcal I_p(0,0,0;s)
=(1-W)^9
\left[1+\big(1-p^{2s-1}\big)
 \sum_{j\ge1}G_4(1;p^j)^2W^j\right].\] The \(j\ge3\) tail is \(O(W^3)\), since the coefficients grow at most polynomially in \(j\), uniformly for large \(p\). Therefore \[\mathcal I_p(0,0,0;-1/10)
=1-9p^{-8/5}+O(p^{-9/5}).\] In contrast, [eq:printed] (\#eq:printed) would give \[\mathcal I_p(0,0,0;-1/10)
=1+9p^{-8/5}+O(p^{-173/100}).\] The difference \(18p^{-8/5}\) cannot be absorbed by \(O(p^{-173/100})\) as \(p\) tends to infinity through primes. This is an analytic specialization within the stated domain, not a floating-point experiment.\par
Impact, adverse checks, and correction. The verified local correction is to replace \(1-Y_pp^{-2-4s}\) by \(1+Y_pp^{-2-4s}\), retaining the stated definition of \(Y_p\). The exact Euler factor is unchanged. The error does not reverse the absolute-convergence argument, because that argument uses the magnitude and cancellation of the first-order coefficient. For example, direct estimates from the exact factor give \[\mathcal I_p-1\ll
p^{-2-4\sigma+4\delta'}
+p^{-2-2\sigma+2\delta'}
+p^{-2+3\delta'}.\] Each exponent is strictly less than \(-1\) when \(\sigma>-1/4+\delta'\) and \(0<\delta'<1/5\). Thus a stronger direct bound supplies the required convergence in that region; the displayed coarse remainder alone need not be used for this inference.\par
The proof’s sign and the leading-constant checks agree with the correction. No central failure of the eighth-moment asymptotic is established from this defect. It is classified Minor because its demonstrated effect is a local coefficient statement whose sign is not needed for the inspected convergence and leading-constant uses. The classification is not based merely on the brevity of the repair.\par
\section*{Checks of the principal proof mechanism}
\subsection*{Truncation with an independent small-parameter constant}
Proposition 3.1, pp. 4184–4186 (PDF 10–12), replaces the natural length parameter by \(U=T^{1-\varepsilon}\). For \(t\asymp T\), \[\left|\frac{t^{4i\tau}-U^{4i\tau}}{i\tau}\right|
\le4|\log(t/U)|\ll\varepsilon\log T\] once \(T\ge T_\varepsilon\). The bounded term \(\log(t/T)\) is absorbed by this permitted threshold; the resulting multiplicative constant need not depend on \(\varepsilon\).\par
The needed shifted bound has exponent \(k=4\) and opposite shifts. The sensitive near-zero integral is \[\int_{|\tau|\le1/200}
 \min\{(2|\tau|)^{-8},(\log T)^8\}\,d\tau
\ll(\log T)^7.\] Combined with \(\varepsilon\log T\) and \(T(\log T)^8\), this gives \(O(\varepsilon T(\log T)^{16})\). Rapid decay controls the distant-shift terms. This explains why an extra \((\log T)^\varepsilon\) loss would have been inadequate.\par
\subsection*{The range and cost of applying Conjecture 1}
The input matching in Section 6.2 and Lemma 6.3, pp. 4200–4201 (PDF 26–27), was checked with the conjecture’s fixed \(C\). The dyadic restrictions imply \[M\asymp N,\quad MN\ll U^{4+\varepsilon_2},\quad
M\ll T^{2-3\varepsilon/2},\] when \(\varepsilon_2\le\varepsilon\). For \(T^{1/2}\le T_0\ll T\) and \(\varepsilon'\le\varepsilon_1<1/6\), \[(MT_0)^{\varepsilon_1}\ll T^{3\varepsilon_1}
=o(T^{1/2})\le T_0.\] Hence the needed shift range \(|r|\ll MT_0^{-1+\varepsilon_1}\) lies within \(|r|\ll M^{1-\varepsilon'}\).\par
The normalization \(T^{-4\varepsilon_1}f_r\) and \(P=T^{1+\varepsilon_1}/T_0\) retain the derivative cost. After counting shifts, \[E_{M,N}\ll
(T/T_0)^{1+C}T^{(5+C)\varepsilon_1}M^{1/2+\varepsilon_0}.\] The stated parameter conditions \(\max(\varepsilon_0,\varepsilon_1,C\varepsilon_1)\le\varepsilon/40\) leave room for \[(5+C)\varepsilon_1+2\varepsilon_0\le\varepsilon/5<\varepsilon/4.\] The dyadic geometric sum of \(M^{1/2}\) is \(O(U^{1+\varepsilon_2/4})=O(T^{1-3\varepsilon/4})\). Thus the claimed total \((T/T_0)^{1+C}T^{1-\varepsilon/2}\) is compatible with every counted loss. Extending the resulting explicit main-term expression in \(r\) does not mean that the conjecture was applied outside its stipulated range.\par
\subsection*{A removable pole}
On p. 4212 (PDF 38), the apparent pole at \(s_*=(z_1+z_2-1)/2\) is cancelled exactly. Put \(a=(z_2-z_1)/2\). The Gamma-ratio sum at \(s_*\) is \[\frac{\Gamma(a+it)}{\Gamma(1+a+it)}
+\frac{\Gamma(-a-it)}{\Gamma(1-a-it)}
=\frac1{a+it}+\frac1{-a-it}=0.\] Its denominators are nonzero for \(t\asymp T\) and the stated small circles. The other Gamma factor is \(\Gamma(1)\). Thus this candidate concern is not an omitted residue.\par
\subsection*{Leading coefficients and removal of smoothing}
The diagonal coefficient in (4.8), p. 4191 (PDF 17), follows directly from the order-17 pole: \[2\,\frac{4^{16}}{2^{16}16!}=\frac{131072}{16!}
=\frac4{638512875}.\] The off-diagonal highest-degree coefficient was additionally checked by an independent, reviewer-written exact rational calculation; its retained program and actual output are provided with the report. The resulting coefficient is \(-107048/16!=-13381/2615348736000\). Thus the assembly in Propositions 3.2–3.3, p. 4187 (PDF 13), gives \[131072-107048=24024.\] The factor \((\log U)^{16}=(1-\varepsilon)^{16}(\log T)^{16}\) changes the leading expression by only \(O(\varepsilon(\log T)^{16})\).\par
For smoothing, fix \(0<\eta<1/2\) and set \(T_0=\eta T\). The remaining error \(O(T\eta^{-1-C})\), divided by \(T(\log T)^{16}\), vanishes as \(T\to\infty\) with \(\eta,\varepsilon\) fixed. Only afterwards are the small parameters sent to zero. The initial omitted dyadic interval is controlled by the RH moment bound and a further fixed parameter \(\eta_0\). The argument on pp. 4187–4189 (PDF 13–15) respects this order of limits; it does not require uniformity for an arbitrary shrinking \(\eta(T)\).\par
\section*{Boundary of the assessment}
Theorem 1.2 states a shifted bound for \(k\ge1/2\). The consulted published companion theorem states \(k\ge1\) with an additional bound on the sum of shifts. The present proof uses \(k=4\) and opposite shifts, which satisfy that version. Simple Hölder interpolation does not supply the sharper small-\(k\) exponent, but later work of Curran allows nonnegative moment exponents and provides a route to the wider bound. This is recorded as a citation-range clarification, not as evidence that the bound is false or that the main proof lacks its required input.\par
The article mentions Maple for involved residues. No author Maple program was identified in the frozen source files or bounded publisher/author-source checks, and none was executed. Code absence is not a mathematical defect. The retained rational check independently tests the leading coefficient only; it does not certify every lower-order term or analytic remainder.\par
No helper is labelled a separate web Pro session or a human reviewer. Backend identity was not independently verified.\par
Overall assessment: one established Minor problem; no further substantive issue established in the declared inspected main-theorem scope. Scientific decision and adjudication remain pending, with no human review or independent acceptance claimed.\par
9 N. Ng, Q. Shen, and P.-J. Wong, The eighth moment of the Riemann zeta function, J. Eur. Math. Soc. 28 (2026), no. 10, 4175–4237. Published online 9 January 2025. https://doi.org/10.4171/JEMS/1568. N. Ng, Q. Shen, and P.-J. Wong, Shifted moments of the Riemann zeta function, Canad. J. Math. 76 (2024), no. 5, 1556–1586. Theorem 1.3 statement consulted in publisher HTML. https://doi.org/10.4153/S0008414X23000548. M. J. Curran, Correlations of the Riemann zeta function, arXiv:2303.10123v2, 15 May 2024. Author version identified separately. https://arxiv.org/abs/2303.10123v2.\par
\end{document}
